\begin{frame}{왜 반경이 정확히 0.2인가 + 스텝당 보상}
  팔이 쭉 펴도 \textbf{0.21 m}까지밖에 닿지 않으므로 --- 13.1절의
  ``도달 가능 영역''에서 보았듯, 목표가 팔의 도달 범위를 넘으면 과업
  자체가 \emph{불가능}. \textbf{$0.2 < 0.21$} 의 설계는 ``거의 항상
  닿을 수 있다''를 보장하면서, 간간이 닿기 어려운 목표가 나와 지나치게
  쉬워지지는 않게.
  {\footnotesize 참고: 반경 $R$ 원판의 균등 지점이 \emph{중심}으로부터
  갖는 평균 거리는 $2R/3 \approx 0.133$ m --- 목표는 원점 부근에
  놓인다는 뜻이지, fingertip(약 0.21 m)과의 평균 거리를 주지는 않음.}
  \vskip0.6em
  스텝당 보상은 두 항의 합 (\texttt{info} 에서 따로 볼 수 있음):
  \[
    r = \underbrace{-1.0 \times \|p_{\text{tip}} - p_{\text{target}}\|}
    _{\text{거리 항}}
    \;-\;
    \underbrace{1.0 \times \|a\|^{2}}_{\text{제어 비용}}
  \]
\end{frame}
